\chapter{Realisation}

\section{Mise en place de l'environnement}
La phase initiale a consiste a valider l'environnement d'execution:
\begin{itemize}[leftmargin=1.2cm]
  \item accessibilite du bastion,
  \item lancement du backend FastAPI,
  \item lancement du frontend React,
  \item mise en place d'un tunnel SSH pour l'API,
  \item configuration du kubeconfig pour l'acces cluster.
\end{itemize}

\section{Implementation du backend}
Le backend expose les endpoints necessaires au pipeline de migration:
\begin{itemize}[leftmargin=1.2cm]
  \item decouverte des VMs,
  \item analyse de compatibilite,
  \item generation du plan,
  \item orchestration de migration,
  \item suivi d'etat des jobs.
\end{itemize}

Une attention particuliere a ete portee a l'authentification JWT afin d'aligner le comportement backend avec le frontend.

\section{Implementation du frontend}
Le frontend a ete transforme d'une interface de demonstration vers une interface de pilotage reelle:
\begin{itemize}[leftmargin=1.2cm]
  \item parcours d'authentification,
  \item affichage de sante API,
  \item lancement Analyze et Plan,
  \item affichage des resultats,
  \item suivi des identifiants de migration,
  \item formulaire de migration reelle.
\end{itemize}

\section{Support des VMs VMware locales}
Un mecanisme local a ete introduit pour les VMs non visibles depuis le bastion:
\begin{itemize}[leftmargin=1.2cm]
  \item lecture locale du fichier \texttt{.vmx},
  \item extraction des caracteristiques utiles,
  \item precheck local de compatibilite,
  \item generation locale d'un plan de base.
\end{itemize}

\section{Gestion des disques VMDK split}
Le pipeline a ete adapte pour prendre en charge les disques VMware segmentes
(\texttt{-s001.vmdk}, \texttt{-s002.vmdk}, etc.), afin d'eviter les echecs de conversion.

\section{Integration OpenShift}
La migration reelle repose sur l'utilisation des outils \texttt{oc}, \texttt{virtctl} et \texttt{qemu-img}.
Des ajustements ont ete necessaires sur le stockage cluster (StorageClass, PVC, HostPathProvisioner) pour stabiliser la chaine de deploiement.

\section{Etat d'avancement et parties restantes}
\textbf{Parties finalisees:}
\begin{itemize}[leftmargin=1.2cm]
  \item architecture generale,
  \item pipeline applicatif principal,
  \item precheck local VMware,
  \item integration initiale OpenShift.
\end{itemize}

\textbf{Parties a completer dans le rapport:}
\begin{itemize}[leftmargin=1.2cm]
  \item captures d'ecran frontend/backend,
  \item exemples de payload API,
  \item valeurs exactes des mesures de performance,
  \item sequence detaillee d'une migration reussie de bout en bout.
\end{itemize}
